% ID: SEL-D0006-OB-E02
% ESTADO: validada. S007, 2026-11-07; tres pasadas del mismo autor y revisión matemática, didáctica y visual.
\textbf{a) Monotonía y extremos relativos.}\par\nopagebreak[4]
El dominio es $\mathbb R$. Derivamos y factorizamos:
\[
f'(x)=x^2-4x+3=(x-1)(x-3).
\]
La derivada es positiva en $(-\infty,1)$ y $(3,\infty)$ y negativa en $(1,3)$. Por tanto, $f$ crece en $(-\infty,1)$ y $(3,\infty)$, y decrece en $(1,3)$.

En 1 pasa de creciente a decreciente: máximo relativo con $f(1)=7/3$. En 3 pasa de decreciente a creciente: mínimo relativo con $f(3)=1$.
\[
\boxed{\text{Máximo relativo }(1,7/3);\qquad\text{mínimo relativo }(3,1).}
\]
No son extremos absolutos en todo $\mathbb R$.

\textbf{b) Curvatura e inflexión.}\par\nopagebreak[4]
$f''(x)=2x-4$. Es negativa para $x<2$ (curvatura hacia abajo, cóncava) y positiva para $x>2$ (hacia arriba, convexa). Cambia de signo en 2 y la función es continua, así que el punto de inflexión es
\[
\boxed{(2,f(2))=(2,5/3).}
\]
Se explicita el signo de la segunda derivada para evitar ambigüedad terminológica sobre concavidad y convexidad.

\textbf{c) Pendiente en x=0.}\par\nopagebreak[4]
La pendiente de la tangente es \boxed{m=f'(0)=3}.

\textbf{d) Integral indefinida.}\par\nopagebreak[4]
Integramos término a término usando la regla de potencias:
\[
\boxed{\int f(x)\,dx=\frac{x^4}{12}-\frac23x^3+\frac32x^2+x+K.}
\]
La constante $K$ representa la familia de primitivas. Comprobación: al derivar el resultado se obtiene $x^3/3-2x^2+3x+1=f(x)$. Además, $f''(1)=-2$ y $f''(3)=2$, coherentes con los extremos calculados.
